\paragraph{Lee--Yang 有限分歧平行族的 pro-\texorpdfstring{$2$}{2} 极限、相干超度量与全息读出阈值}
沿用推论 \ref{cor:xi-terminal-zm-leyang-ramification-sum-zero} 与推论 \ref{cor:xi-terminal-zm-leyang-ramification-pullback-2n} 的记号：
设 \(E:=E_{\min}\)，\(\pi_y:E\to \PP^1_y\) 为由有理函数 \(y\) 给出的次数 \(4\) 覆盖；
其五个有限简单分歧点记为 \(P_0,\dots,P_4\)，满足
\[
R(y)=3O+\sum_{i=0}^{4}P_i,
\qquad
\sum_{i=0}^{4}P_i=O.
\]

\begin{theorem}[有限分歧平行族的 pro-\texorpdfstring{$2$}{2} 扭子极限]\label{thm:xi-leyang-pro2-torsor-limit-five-families}
对每个 \(n\ge 0\) 置 \(y_n:=y\circ[2^n]\)。
记其有限分歧集合（忽略映到 \(\infty\) 的全分歧部分）为 \(R_{\mathrm{fin}}(y_n)\)。
则
\[
R_{\mathrm{fin}}(y_n)=\bigsqcup_{i=0}^{4}\bigl(P_i+E[2^n]\bigr).
\]
对每个 \(i\in\{0,\dots,4\}\) 与 \(n\ge0\)，定义
\[
B_{i,n}:=P_i+E[2^n],
\]
以及层间映射
\[
\pi^{(i)}_{n+1,n}:B_{i,n+1}\longrightarrow B_{i,n},
\qquad
\pi^{(i)}_{n+1,n}(P_i+T):=P_i+2T.
\]
记 \(T_2(E):=\varprojlim_n E[2^n]\)。
则逆极限
\[
B_{i,\infty}:=\varprojlim_n(B_{i,n},\pi^{(i)}_{n+1,n})
\]
天然是一个 \(T_2(E)\)-主齐次空间，并且
\[
\Psi_i:T_2(E)\xrightarrow{\ \sim\ }B_{i,\infty},
\qquad
(T_n)_{n\ge0}\longmapsto (P_i+T_n)_{n\ge0}
\]
为显式同构。
因而
\[
\varprojlim_n R_{\mathrm{fin}}(y_n)
\cong
\bigsqcup_{i=0}^{4}B_{i,\infty}
\cong
\bigsqcup_{i=0}^{4}T_2(E).
\]
\end{theorem}
\begin{proof}
由推论 \ref{cor:xi-terminal-zm-leyang-ramification-pullback-2n}，
\[
R(y_n)
=
3\sum_{Q\in E[2^n]}[Q]
+\sum_{Q\in E[2^n]}\sum_{i=0}^{4}[P_i+Q].
\]
故有限分歧支持恰为 \(\bigsqcup_{i=0}^{4}(P_i+E[2^n])\)，即上式中的 \(R_{\mathrm{fin}}(y_n)\)。

对 \(\pi^{(i)}_{n+1,n}\) 的良定义性：若 \(T\in E[2^{n+1}]\)，则 \(2T\in E[2^n]\)，故
\(
P_i+2T\in B_{i,n}
\)。
再看 \(\Psi_i\)：若 \((T_n)_n\in T_2(E)\)，即 \(2T_{n+1}=T_n\)，则
\[
\pi^{(i)}_{n+1,n}(P_i+T_{n+1})=P_i+2T_{n+1}=P_i+T_n,
\]
故 \((P_i+T_n)_n\) 为兼容族。

反之，任给 \((X_n)_n\in B_{i,\infty}\)，可唯一写为 \(X_n=P_i+T_n\)（\(T_n\in E[2^n]\)）。
兼容性给出
\[
X_n=\pi^{(i)}_{n+1,n}(X_{n+1})
\Longleftrightarrow
T_n=2T_{n+1},
\]
故 \((T_n)_n\in T_2(E)\)，且 \(\Psi_i((T_n)_n)=(X_n)_n\)。
于是 \(\Psi_i\) 双射，结论成立。证毕。
\end{proof}

\begin{theorem}[相干时间的非阿基米德化与完备性]\label{thm:xi-coherence-ultrametric-on-pro2-branch-torsor}
固定 \(i\in\{0,\dots,4\}\)，在 \(B_{i,\infty}\) 上定义
\[
\nu(x,y):=\sup\{n\ge0:\ x_n=y_n\}\in\NN\cup\{\infty\},
\]
\[
d_{\mathrm{coh}}(x,y):=
\begin{cases}
0,&x=y,\\
2^{-\nu(x,y)},&x\neq y.
\end{cases}
\]
则 \((B_{i,\infty},d_{\mathrm{coh}})\) 是完备超度量空间。
在同构 \(\Psi_i^{-1}:B_{i,\infty}\to T_2(E)\) 下，若记
\[
m(T,S):=\sup\{n\ge0:\ T-S\in 2^nT_2(E)\},
\]
则
\[
\nu\bigl(\Psi_i(T),\Psi_i(S)\bigr)=m(T,S),
\]
从而 \(d_{\mathrm{coh}}\) 与 \(T_2(E)\) 的标准 \(2\)-进距离等距同构。
\end{theorem}
\begin{proof}
令 \(x=\Psi_i(T)\)、\(y=\Psi_i(S)\)。
由 \(\Psi_i\) 定义有
\[
x_n=y_n
\Longleftrightarrow
P_i+T_n=P_i+S_n
\Longleftrightarrow
T_n=S_n.
\]
而 \(T_n=S_n\) 等价于 \(T-S\in 2^nT_2(E)\)，故
\(
\nu(\Psi_i(T),\Psi_i(S))=m(T,S)
\)。

超度量不等式由 \(2\)-进滤过立即给出：若
\(
T-S\in 2^aT_2(E)
\) 且
\(
S-U\in 2^bT_2(E)
\)，则
\[
T-U=(T-S)+(S-U)\in 2^{\min(a,b)}T_2(E),
\]
故
\[
\nu(x,z)\ge\min\{\nu(x,y),\nu(y,z)\},
\qquad
d_{\mathrm{coh}}(x,z)\le\max\{d_{\mathrm{coh}}(x,y),d_{\mathrm{coh}}(y,z)\}.
\]
完备性来自 \(T_2(E)\) 的 \(2\)-进完备性与 \(\Psi_i\) 的等距性。证毕。
\end{proof}

\begin{theorem}[Gödel--Tate 块编码的等距同余转写]\label{thm:xi-godel-tate-block-coding-isometry}
固定 \(T_2(E)\) 的一个 \(\ZZ_2\)-基 \((e_1,e_2)\)。
则任意 \(T\in T_2(E)\) 唯一写为
\[
T=\sum_{k\ge0}\bigl(a_k2^k e_1+b_k2^k e_2\bigr),
\qquad
a_k,b_k\in\{0,1\}.
\]
定义块式编码
\[
\Gamma:T_2(E)\to\ZZ_2,
\qquad
\Gamma(T):=\sum_{k\ge0}(a_k+2b_k)\,4^k.
\]
则 \(\Gamma\) 为双射，且对任意 \(T,S\in T_2(E)\) 与 \(n\ge0\) 有
\[
T\equiv S\ (\mathrm{mod}\ 2^nT_2(E))
\Longleftrightarrow
\Gamma(T)\equiv \Gamma(S)\ (\mathrm{mod}\ 4^n\ZZ_2).
\]
进一步，对任意固定分支 \(i\) 与 \(T\neq S\)（约定 \(v_2(0)=+\infty\)）,
\[
\nu\bigl(\Psi_i(T),\Psi_i(S)\bigr)
=
\left\lfloor\frac{v_2\!\bigl(\Gamma(T)-\Gamma(S)\bigr)}{2}\right\rfloor.
\]
\end{theorem}
\begin{proof}
\((a_k,b_k)\) 与 \((a_k',b_k')\) 分别是 \(T,S\) 的二进系数对。
同余 \(T\equiv S\ (\mathrm{mod}\ 2^nT_2(E))\) 当且仅当前 \(n\) 级系数对完全一致，即
\[
(a_k,b_k)=(a_k',b_k')\quad (0\le k<n).
\]
这等价于 \(\Gamma(T)\) 与 \(\Gamma(S)\) 的前 \(n\) 个 \(4\)-进位一致，即
\[
\Gamma(T)\equiv\Gamma(S)\ (\mathrm{mod}\ 4^n\ZZ_2).
\]
故同余等价式成立，且 \(\Gamma\) 双射（每个 \(4\)-进位 \(0,1,2,3\) 与二比特块一一对应）。

再由
\[
\Gamma(T)-\Gamma(S)\equiv0\ (\mathrm{mod}\ 4^n)
\Longleftrightarrow
v_2\!\bigl(\Gamma(T)-\Gamma(S)\bigr)\ge 2n
\]
与定理 \ref{thm:xi-coherence-ultrametric-on-pro2-branch-torsor} 的 \(\nu=m\) 识别，得到估值公式。证毕。
\end{proof}

\begin{theorem}[三纤维张量化的六维 \texorpdfstring{$2$}{2}-进超立方体逆极限]\label{thm:xi-three-fiber-tensorized-sixdim-pro2-hypercube}
取三份独立的 Lee--Yang 分支扭子读出，固定分支指标 \(i,j,k\in\{0,\dots,4\}\)，定义
\[
\mathcal B_\infty
:=
B_{i,\infty}^{(1)}\times B_{j,\infty}^{(2)}\times B_{k,\infty}^{(3)}
\cong
T_2(E)^3
\cong
\ZZ_2^6.
\]
对每层 \(n\ge0\)，其截断层
\[
\mathcal B_n
:=
B_{i,n}^{(1)}\times B_{j,n}^{(2)}\times B_{k,n}^{(3)}
\cong
E[2^n]^3
\cong
(\ZZ/2^n\ZZ)^6,
\]
故
\[
|\mathcal B_n|=2^{6n}.
\]
在选定 \(\ZZ_2\)-基后，可把 \(\mathcal B_n\) 识别为 \(6n\)-比特串集合 \(\{0,1\}^{6n}\)（等价于超立方体 \(Q_{6n}\) 的顶点集）；
并且层间投影 \(\mathcal B_{n+1}\to\mathcal B_n\) 对应“去最高位块截断”（每步删去 \(6\) 比特）。
\end{theorem}
\begin{proof}
由定理 \ref{thm:xi-leyang-pro2-torsor-limit-five-families}，每个分支层 \(B_{i,n}^{(\ell)}\) 与 \(E[2^n]\) 同构；
选定基后 \(E[2^n]\cong(\ZZ/2^n\ZZ)^2\)，三份直积即得 \((\ZZ/2^n\ZZ)^6\)。
每个坐标有 \(n\) 个二进位，故与 \(6n\)-比特串一一对应，基数即 \(2^{6n}\)。
层间映射由模 \(2^n\) 约化给出，正是各坐标删去最高位；六坐标并置后即“每层删去 \(6\) 比特块”。证毕。
\end{proof}

\begin{theorem}[全息读出与相干深度的指数互易律]\label{thm:xi-holographic-readout-coherence-exponential-law}
令 \(G:=\ZZ_2^d\)（\(d\ge1\)），取 \(0\le m\le n\)。
令随机元 \(A\) 在商群 \(G/2^nG\) 上均匀分布，记自然投影为
\[
\pi_m:G/2^nG\to G/2^mG.
\]
设 \(\mathsf{Alg}\) 为任意（可随机化）决策算法，其输入仅为 \(\pi_m(A)\) 与独立随机源，输出候选 \(C\in G/2^nG\)。
则几乎处处有
\[
\PP\!\bigl[C=A\mid \pi_m(A)\bigr]\le 2^{-d(n-m)}.
\]
因此若希望成功概率至少为 \(\eta\in(0,1)\)，则必有
\[
m\ge n-\frac1d\log_2\frac1\eta.
\]
在独立盲试制度下，达到一次命中的期望试验次数至少为
\[
2^{d(n-m)}.
\]
特别地，\(d=2\) 与 \(d=6\) 分别给出 \(4^{\,n-m}\) 与 \(64^{\,n-m}\) 级别壁垒。
\end{theorem}
\begin{proof}
条件在 \(\pi_m(A)=a_m\) 下，\(A\) 在余类 \(a_m+2^mG/2^nG\) 上均匀分布。
该余类基数为
\[
|2^mG/2^nG|=2^{d(n-m)}.
\]
算法输出 \(C\) 在该条件下至多指定其中一个元素，故条件命中概率不超过其倒数，得到第一式。
其余结论由不等式代数变形与几何分布期望下界 \(1/p\)（\(p\le 2^{-d(n-m)}\)）直接推出。证毕。
\end{proof}

\begin{proposition}[Zeckendorf 归一化携带深度的对数局域性]\label{prop:xi-zeckendorf-fold-carry-depth-log-locality}
令 \(\omega\sim\mathrm{Unif}(\{0,1\}^m)\)，并令 \(\Fold_m(\omega)\) 表示 Fibonacci 权重下的 Zeckendorf 归一化。
记 \(D(\omega)\) 为归一化过程中携带链在索引轴上的最大传播跨度。
则存在与 \(m\) 无关的常数 \(c_1,c_2>0\)，使任意 \(L\ge1\) 满足
\[
\PP\!\bigl[D(\omega)\ge L\bigr]\le c_1\,m\,2^{-L},
\]
并且
\[
\EE[D(\omega)]\le c_2\log_2 m.
\]
\end{proposition}
\begin{proof}
在标准局部重写规则下，若发生长度至少 \(L\) 的携带传播链，则传播路径上必须连续触发至少 \(L\) 次局部相邻约束；
这要求输入字中存在长度至少 \(L\) 的连续 \(1\)-块。
故有集合包含
\[
\{D(\omega)\ge L\}
\subseteq
\{\omega\ \text{含长度}\ \ge L\ \text{的连续}\ 1\text{-块}\}.
\]
后者以并合估计得
\[
\PP[D(\omega)\ge L]
\le (m-L+1)\,2^{-L}
\le m\,2^{-L},
\]
即第一式（取 \(c_1=1\) 即可）。

再用尾和公式
\[
\EE[D]=\sum_{L\ge1}\PP[D\ge L]
\]
并按 \(L\le\lfloor\log_2 m\rfloor\) 与 \(L>\lfloor\log_2 m\rfloor\) 分段：
\[
\EE[D]
\le
\sum_{L\le\lfloor\log_2 m\rfloor}1
+\sum_{L>\lfloor\log_2 m\rfloor}m2^{-L}
\le
\log_2 m+2.
\]
因此存在常数 \(c_2\)（例如 \(c_2=3\)）使所述上界成立。证毕。
\end{proof}

\begin{conjecture}[全息--相干公设的算术刚性与 conductor \texorpdfstring{$37$}{37} 唯一性]\label{conj:xi-holography-coherence-arithmetic-rigidity-conductor37}
设某扫描--投影生成器 \(\mathscr G\) 在全部分辨率层同时满足以下三项结构条件：
\begin{enumerate}
  \item 每层有限分歧集合可精确分解为五个互不相交的平行族，并且与精度截断投影形成可测因子系统；
  \item 存在有限数字集驱动的 pro-\(2\) 仿射 IFS 模型，使“时间推进 = 精度提升”以及“相干 = 进入同一 \(2^n\)-球”的对应在 \(2\)-进超度量下严格成立；
  \item 分歧判别式通道与 Lee--Yang 三次 \(P_{\mathrm{LY}}\) 的算术分歧素集合对齐。
\end{enumerate}
则 \(\mathscr G\) 在可测共轭意义下与 \(T_2(E_{\min})\) 上的仿射系统同型；
并且椭圆锚点在 \(\QQ\)-同构意义下由 \(\mathrm{cond}(E_{\min})=37\) 唯一确定。
\end{conjecture}

\endinput
